\scanpage{7}
% L1–2
topos $E,F$. On suppose comme dans la dém. de prop 1 $E=\widetilde C$, $F=\widetilde D$, \deleted{\unsure{donc}} on a vu que l'on a alors

% L3–4
\[
\begin{aligned}
\pi(E,F)=\pi(\unsure{\overline C|\widetilde C},\widetilde D)
 &=\deleted{\unsure{\pi}}\pi(\widetilde C,\widehat D)\cap\Pi(\widehat C,\widetilde D)\\
 &\quad\text{(intersection des sous topos de }\pi(\widehat C,\widehat D)\simeq\widehat{C\times D}\text{)}
\end{aligned}
\]

% L5–9
et il est clair que les deux facteurs $\pi(\widetilde C,\widehat D)$ et $\pi(\widehat C,\widetilde D)$ sont les images inverses des sous topos $\widetilde C\subset\widehat C$ resp. $\widetilde D\subset\widehat D$ par les ``projections'' $\Pi(\widehat C,\widehat D)\longrightarrow\widehat C$ resp. $\pi(\widehat C,\widehat D)\longrightarrow\widehat D$. Cela nous ramène \unsure{à|au} \deleted{cas où} prouver que $\pi(\unsure{C|\widehat C},\unsure{\overline D|\widehat D})$ est un 2-produit des topos $\widehat C,\widehat D$. Or pour un topos \deleted{E} $X$, on a

% L10–11
\[
\begin{aligned}
\underline{\mathrm{Hom}}_{\mathrm{top}}(X,\widehat C)&\cong\underline{\mathrm{Hom}}'(C,X)^o\\
\underline{\mathrm{Hom}}_{\mathrm{top}}(X,\widehat D)&\cong\underline{\mathrm{Hom}}'(D,X)^o
\end{aligned}
\]

% L12–13
où le $\underline{\mathrm{Hom}}'$ désigne les foncteurs exacts à g [NB on a supposé $C$ et $D$ stables par $\varprojlim$ finies], donc on a

% L14–16
\[
\begin{aligned}
\underline{\mathrm{Hom}}_{\mathrm{top}}(X,\widehat C)\times\underline{\mathrm{Hom}}_{\mathrm{top}}(X,\widehat D)
 &\cong[\underline{\mathrm{Hom}}'(C,X)\times\underline{\mathrm{Hom}}'(D,X)]^o\\
 &\cong(\underline{\mathrm{Hom}}'(C\times D,X))^o\\
 &\cong\underline{\mathrm{Hom}}_{\mathrm{top}}(X,\unsure{\widehat{C\times D}|\widehat C\times D}))\qquad\text{cqfd}
\end{aligned}
\]

% L17–21
\inserted{[si $C$ et $D$ sont stables par $\varprojlim$ finies, \deleted{il} \deleted{\unsure{C}} en est de $=$ de $C\times D$, et les foncteurs exacts à gauche \inserted{(Attention, en les couples de variables, non pas en les variables séparément)} $C\times D\xrightarrow{w}X$ [$X$ une catégorie quel.] correspondent aux \deleted{\unsure{fon|deux}} couples de foncteurs exacts à gauche $C\xrightarrow{u}X$, $D\xrightarrow{v}X$, en leur associant \deleted{\unsure{$u(a)$}} $(a,b)\longrightarrow u(a)\times v(b)$, et à $w$ les foncteurs $a\longrightarrow w(a,e_D)$, $b\longrightarrow w(e_C,b\unsure{)}$}

% foot
% archive sheet number: 6
